\chapter{Conclusiones del proyecto}
\section {Conclusiones de los resultados obtenidos}
Teniendo en cuenta el trabajo de investigación llevado a cabo y los resultados finales obtenidos del análisis de seguridad, se pueden extraer las siguientes conclusiones:
\begin{itemize}
\item Las vulnerabilidades encontradas y explotadas en este proyecto, son el resultado de errores de diseño o falta de implementación de mecanismos de seguridad en los dispositivos y redes IoT. Con esto, se ha puesto de manifiesto la importancia que la seguridad debe tener cuando se elige una tecnología de comunicaciones inalámbrica, para interconectar cualquier par de dispositivos, y, sobre todo si estos se utilizan en el hogar recopilando datos sensibles.
\item Los mecanismos de seguridad proporcionados por los protocolos Thread y Zigbee se pueden considerar adecuados y robustos, ya que se utilizan estándares como AES para el cifrado y la autenticación de los datos y se disponen de mecanismos de distribución segura de las claves y de protección ante duplicaciones. No obstante, esta seguridad depende de la implementación y configuración que realicen los desarrolladores sobre el protocolo.
\item Aunque Thread y Zigbee estén basados en el mimso estándar, al hacer una comparativa entre las medidas de seguridad establecidas por cada protocolo, se comprueba que Thread ofrece mayor robustez y seguridad debido principalmente a tres factores:
	\begin{itemize}
	\item La seguridad a nivel de red es obligatoria, a contraposición con Zigbee donde es opcional.
	\item Únicamente se pueden unir dispositivos a la red que cuenten con una clave preinstalada mediante un método seguro.
	\item La gran cantidad de tipos de dispositivos Thread en una topología en malla no permite que haya un único punto de fallo.
	\end{itemize}
\end{itemize}
\clearpage



\section {Futuras líneas de investigación}
A continuación se indican una serie de líneas de investigación que en base a este proyecto se podrían desarrollar en futuros trabajos:
\begin{itemize}
\item Efectuar una auditoría de seguridad en una red Thread.
\item Realizar un estudio de la seguridad física en los dispositivos IoT utilizados en este proyecto (fabricante Ksix).
\item Analizar las comunicaciones y configuraciones por defecto, orientadas a la seguridad de la red, de otros fabricantes de dispositivos IoT como Xiaomi o Philips.
\item Ejecutar ataques contra una red IoT que necesiten ser efectuados desde varios puntos, disponiendo para ello de dos o más dispositivos Api-Mote.
\item Ejecutar ataques contra una red IoT que disponga de medidas de seguridad más efectivas (sistema de detección de intrusos) y una infraestructura mayor (varios enrutadores para redirigir el tráfico de red).
\item Realizar un estudio de ahorro energético y económico al construir un edificio o casa inteligente, incluyendo una red de dispositivos IoT.
\end{itemize}



